\chapter{À leur manière}
\label{chap:a-leur-maniere}

Les années passées sur les routes avaient offert à Khaer Ghal et Kamatsu bien
plus que les leçons données par les voyageurs qu'ils rencontraient. Entre les
corvées, les haltes et les longues journées de marche, les deux garçons avaient
aussi appris à occuper leur temps comme tous les enfants de leur âge : en
inventant des jeux, en explorant les environs et, parfois, en faisant des
choses qu'ils savaient pertinemment ne pas devoir faire.

Un jour, ils décidèrent de fabriquer leurs propres arcs.

L'idée leur était venue après avoir passé plusieurs jours à observer un
chasseur qui voyageait avec la caravane. Ils l'avaient vu entretenir son arme,
changer une corde et examiner différents morceaux de bois en leur expliquant
ce qui pouvait faire une bonne branche. Lorsqu'il était reparti, les deux
garçons avaient retenu suffisamment de choses pour être convaincus qu'ils
étaient désormais capables de fabriquer les leurs.

La réalité se révéla un peu plus compliquée.

Le premier morceau de bois choisi par Kamatsu était trop rigide. Celui de
Khaer Ghal se courba correctement jusqu'au moment où il força un peu trop et
le brisa. Ils recommencèrent, comparèrent leurs essais, modifièrent ce qui leur
semblait mauvais.

Après plusieurs tentatives, un arc finit malgré tout par prendre forme.

Ils étaient loin d'avoir fabriqué une arme remarquable. Un véritable archer
aurait sans doute trouvé beaucoup à redire sur leur travail, mais l'arc lançait
des flèches et cela suffisait largement aux deux garçons.

Kamatsu tira sur la corde et observa sa courbure.

— Kael, regarde.

Khaer Ghal leva les yeux.

— Quoi ?

— Je crois qu'on a réussi.

Il prit l'arc que lui tendait son ami et l'examina.

Cela faisait déjà quelque temps que Kamatsu avait commencé à l'appeler ainsi.
Khaer Ghal était trop long à son goût, et les premières tentatives pour le
raccourcir avaient fini par donner Kael. Khaer Ghal l'avait corrigé plusieurs
fois au début, avant de comprendre que cela ne servait absolument à rien.
Kamatsu avait décidé que ce serait Kael et, entre eux, le surnom était
simplement resté.

Pour tous les autres, il était encore Khaer Ghal, ou l'une des nombreuses
variantes plus ou moins heureuses que les voyageurs parvenaient à prononcer.

Pour Kamatsu, il était Kael.

Il rendit l'arc à son ami.

— Essaie.

Ils passèrent une bonne partie de l'après-midi à se partager l'arc et à tirer
sur un morceau de bois posé contre un arbre. La plupart de leurs premières
flèches manquèrent leur cible. Certaines partirent si loin qu'ils passèrent
davantage de temps à les chercher qu'à tirer, mais cela ne les empêcha pas de
recommencer jusqu'à ce que la lumière commence à décliner.

Lorsqu'ils rentrèrent au campement, Kamatsu portait leurs quelques flèches et
Khaer Ghal leur arc.

Ils avaient trouvé une nouvelle manière de s'occuper.

\scenebreak

Quelques jours plus tard, Kael partit seul.

Il avait emporté l'arc et quelques flèches avec l'intention de s'entraîner
pendant que la caravane profitait d'une halte suffisamment longue pour
effectuer plusieurs réparations. Il aurait parfaitement pu rester à proximité
des chariots, mais les arbres qui entouraient le campement lui offrirent
rapidement quelque chose de plus intéressant.

Une piste traversait les sous-bois.

Kael la suivit.

Il ne savait pas vraiment à quel animal appartenaient les empreintes et voulait
le découvrir. La piste finit par disparaître sur un terrain plus dur, mais il
avait alors aperçu une hauteur entre les arbres et décida d'aller voir ce qui
se trouvait au sommet.

De là, il distingua un ruisseau.

Il descendit jusqu'à lui.

Puis il continua.

Il n'avait aucune raison particulière d'aller plus loin. Il ne cherchait rien
et personne ne lui avait confié de tâche. Chaque fois qu'il atteignait
l'endroit qu'il avait voulu voir, quelque chose d'autre attirait simplement
son attention un peu plus loin.

Alors il allait voir.

Il tira quelques flèches, suivit des traces qu'il ne connaissait pas, escalada
un rocher pour observer les environs et continua ainsi sans véritablement
s'interroger sur le temps qui passait.

Ce fut seulement lorsque la lumière devint franchement mauvaise qu'il regarda
le ciel entre les branches.

Le soleil avait disparu.

Kael resta immobile quelques secondes.

— Merde.

Il reprit immédiatement le chemin de la caravane.

Lorsqu'il aperçut enfin les premières lueurs du campement entre les arbres, la
nuit était tombée depuis longtemps. Plusieurs lanternes avaient été allumées
et des voix résonnaient à la lisière de la forêt.

Quelqu'un le cherchait.

Il accéléra.

Kamatsu fut le premier à l'apercevoir.

— Kael !

Il courut vers lui. Le soulagement apparut immédiatement sur son visage, avant
de laisser presque aussitôt place à la colère.

— T'étais où ?

— Dans la forêt.

— Ça, j'avais compris !

D'autres caravaniers approchaient déjà. Kael comprit rapidement, au nombre de
personnes parties à sa recherche et à l'expression de celles restées au camp,
que son absence avait provoqué bien plus d'inquiétude qu'il ne l'avait
imaginé.

Il eut droit à plusieurs remarques, puis à une véritable réprimande. On lui
expliqua tout ce qui aurait pu lui arriver, qu'il aurait dû prévenir quelqu'un
avant de partir et que personne ne savait où commencer à le chercher puisqu'il
n'avait même pas indiqué dans quelle direction il comptait aller.

Kael ne protesta pas.

Plus on lui parlait, plus il regardait ses pieds. Il n'avait pas voulu
inquiéter qui que ce soit et n'avait même pas pensé que quelqu'un pourrait
s'inquiéter. Il avait seulement pris son arc et suivi ce qu'il avait envie de
découvrir, sans demander à personne s'il pouvait partir ni jusqu'où il avait le
droit d'aller.

Lorsqu'on le laissa enfin rejoindre Kamatsu, celui-ci resta silencieux quelques
secondes.

— T'es vraiment idiot.

— Je sais.

— On savait même pas où t'étais.

Kael releva légèrement les yeux vers lui.

— Je sais.

— J'ai cru qu'il t'était arrivé quelque chose.

Cette fois, Kael ne répondit pas immédiatement. La colère de Kamatsu avait
presque disparu et il comprit qu'elle cachait surtout la peur qu'il avait eue
pendant son absence.

— Désolé.

Kamatsu le regarda encore un instant, puis désigna l'arc qu'il tenait toujours.

— Au moins, ça valait le coup ?

Kael hésita.

Il repensa au ruisseau, aux traces qu'il n'avait pas réussi à identifier et au
paysage qu'il avait découvert depuis le sommet du rocher.

Un petit sourire apparut malgré lui.

— Oui.

Kamatsu leva les yeux au ciel.

Kael regrettait sincèrement d'avoir inquiété tout le monde.

Il ne regrettait pas vraiment d'être parti.

\scenebreak

Ce ne fut évidemment pas leur dernière mauvaise idée.

Au cours d'une autre halte, les deux garçons s'étaient éloignés du campement
avec leur arc. Bartiméus les avait suivis, comme il lui arrivait parfois de le
faire lorsqu'ils ne s'éloignaient pas trop, tantôt marchant quelques mètres
devant eux, tantôt disparaissant dans les broussailles avant de réapparaître
un peu plus loin.

Ils finirent par atteindre un petit ravin qu'un arbre tombé permettait de
traverser. Le passage n'avait rien d'impressionnant : le fond se trouvait à
peine quelques mètres plus bas et il était facile de le contourner.

Bartiméus, lui, ne s'encombra pas de telles considérations.

Le chat roux sauta sur le tronc et s'y engagea tranquillement. Il avançait avec
l'assurance naturelle d'un animal qui ne semblait même pas comprendre que
tenir en équilibre puisse représenter une difficulté. Quelques instants plus
tard, il atteignit l'autre côté, sauta au sol et poursuivit son exploration
comme si rien de remarquable ne venait de se produire.

Kamatsu suivit le chat du regard.

Puis il regarda le tronc.

Puis Kael.

— Quoi ? demanda celui-ci.

Kamatsu eut ce sourire que son ami commençait à connaître suffisamment pour
savoir qu'une mauvaise idée venait probablement de lui traverser l'esprit.

— Rien.

— Quoi ?

— Bartiméus vient de traverser.

Kael regarda le chat roux qui s'éloignait déjà de l'autre côté.

— Oui.

— Je parie que t'es même pas capable de faire pareil.

Kael reporta son attention sur le tronc.

— Si.

— Sans les mains ?

La surface était irrégulière, encore couverte par endroits de mousse et
d'écorce humide. Bartiméus l'avait traversée sans la moindre hésitation.

— Si.

Kamatsu sourit davantage.

— Cap ?

Kael aurait pu contourner le ravin. Il aurait pu répondre qu'il n'était pas un
chat et que comparer leur équilibre n'avait aucun sens. Il aurait même pu
rappeler à Kamatsu qu'ils étaient venus ici pour essayer leur arc et non pour
imiter Bartiméus.

À la place, il posa l'arc contre un rocher.

— Cap.

Il grimpa sur le tronc.

Les premiers pas furent faciles. Kael écarta légèrement les bras pour conserver
son équilibre et avança en regardant devant lui. De l'autre côté, Bartiméus
s'était arrêté et semblait maintenant observer la scène. Kamatsu, lui, avait
cessé de sourire et suivait chacun des mouvements de son ami avec beaucoup
plus d'attention.

Kael atteignit la moitié, puis les deux tiers.

— Tu vois ? lança-t-il.

— T'es pas encore arrivé !

Kael continua.

Il ne lui restait que quelques pas lorsqu'il posa le pied sur une plaque de
mousse humide. Sa chaussure glissa et il essaya de retrouver son équilibre,
agitant les bras avant de parvenir presque à se rattraper.

Presque.

Il bascula sur le côté et disparut du tronc.

— Kael !

Il heurta le sol en contrebas dans un bruit de feuilles et de branches.

Kamatsu dévala aussitôt la pente.

— Ça va ?

Kael était assis au milieu des feuilles, le visage crispé.

— Oui.

— T'es sûr ?

— Oui.

Il posa une main au sol et essaya de se relever. Une douleur brutale traversa
sa cheville et il retomba immédiatement.

— Non.

Kamatsu le regarda. Kael regarda sa cheville, puis tous les deux levèrent les
yeux vers le tronc au-dessus d'eux.

Bartiméus était toujours là.

Le chat les observait depuis l'autre côté, parfaitement immobile.

— C'était complètement débile, dit Kamatsu.

Kael lui lança un regard.

— C'est toi qui m'as défié.

— Et toi, t'étais pas obligé de le faire !

Kael ne trouva rien à répondre à cela.

Kamatsu passa finalement son bras autour de ses épaules pour l'aider à se
relever. Sa cheville n'avait rien de grave, mais elle était suffisamment
douloureuse pour qu'il soit incapable de marcher normalement.

Ils reprirent donc lentement la direction du campement, Kael appuyé contre son
ami. Bartiméus finit par les rejoindre et trottina tranquillement à leurs
côtés, manifestement beaucoup moins impressionné par l'aventure que les deux
garçons.

Après quelques minutes, Kamatsu rompit le silence.

— T'étais presque arrivé.

Kael tourna la tête vers lui.

— Je sais.

— Presque.

— Tais-toi.

Kamatsu éclata de rire.

Kael essaya de lui donner un coup d'épaule, ce qui lui rappela immédiatement
qu'il avait besoin de lui pour marcher.

Cela ne fit rire Kamatsu que davantage.

\scenebreak

Les semaines continuèrent ainsi, entre les routes, les apprentissages, les
corvées et les petites aventures que les deux garçons trouvaient toujours le
moyen de transformer en histoires à raconter.

Ils grandissaient ensemble, apprenaient ensemble et faisaient parfois ensemble
des choses suffisamment stupides pour se demander ensuite lequel des deux en
avait eu l'idée.

Kael continuait d'observer tout ce qu'il pouvait, de s'entraîner avec son arc
et de profiter des haltes pour découvrir ce qui se trouvait autour de lui.
Parfois avec Kamatsu, parfois seul.

Et Kamatsu avait découvert quelque chose d'autre.

Lorsqu'il voulait être absolument certain que son ami essaierait quelque
chose, il existait une méthode particulièrement efficace.

Il suffisait de lui dire qu'il n'en était pas capable.